% ============================================================
%  ppo_queueing_report.tex
%  Technical note: PPO for the deadline-constrained multi-server
%  queue --- the method, the stabilization recipe, and benchmark
%  results on Experiments 2 and 3.
%
%  Intended reader: comfortable with MDPs, uniformisation and
%  RVI; no prior exposure to policy-gradient RL assumed.
% ============================================================
\documentclass[11pt,a4paper]{article}

\usepackage{amsmath,amssymb,amsthm}
\usepackage{geometry}
\geometry{margin=2.5cm}
\usepackage{booktabs}
\usepackage{array}
\usepackage{caption}
\usepackage{hyperref}
\usepackage{xcolor}
\usepackage{microtype}
\usepackage{parskip}

\newtheorem{remark}{Remark}
\newtheorem{mechanism}{Mechanism}

\newcommand{\E}{\mathbb{E}}
\newcommand{\R}{\mathbb{R}}
\newcommand{\clip}{\operatorname{clip}}

\title{Proximal Policy Optimization for a Deadline-Constrained
       Multi-Server Queue:\\
       Method, Stabilization Recipe, and Benchmark Results}
\author{Technical Note}
\date{}

\begin{document}
\maketitle

\begin{abstract}
We describe how Proximal Policy Optimization (PPO) was made to learn
(near-)optimal scheduling policies for our deadline-constrained
multi-server queueing MDP, benchmarked against exact Relative Value
Iteration (RVI) solutions.
Section~\ref{sec:ppo} introduces PPO from first principles for a reader
familiar with MDPs and RVI but not with policy-gradient reinforcement
learning.
Section~\ref{sec:ours} presents our specific formulation: the semi-MDP
discounting correction, potential-based reward shaping, value-target
normalization, staggered environment resets, and guarded temperature
annealing --- each introduced as the fix to a specific, diagnosed failure
mechanism.
Section~\ref{sec:experiments} defines the two benchmark instances
(Experiment~3: routing with a specialist and a generalist server;
Experiment~2: fully flexible servers with asymmetric costs, where the
optimal policy requires \emph{strategic idling}) and reports results:
Experiment~3 is solved (7--8/8 seeds within 5\% of RVI), Experiment~2 is
solved per problem instance (7/8 seeds within 6\%; failures are a seed
lottery mitigated by training 2--3 seeds and selecting by evaluation).
A controlled single-ingredient ablation study (8 cells $\times$ 8 seeds)
substantiates each ingredient of the recipe: removing any one of the four
main stabilizers collapses the success rate from 6--7/8 to 0--3/8, with
failures landing on two exactly reproducible attractors (never-serve and
FIFO-clone).
\end{abstract}

\tableofcontents
\newpage

% ============================================================
\section{Setting (recapitulation)}
\label{sec:setting}
% ============================================================

We briefly restate the system and its MDP; full details are in the
companion RVI note.

\paragraph{System.}
$K$ server pools serve $N$ job types.  Jobs of type $n$ arrive as
independent Poisson streams with rate $\lambda_n$ and wait in per-type
FIFO queues.  Service is exponential; server pool $k$ can serve a subset
of types at rates $\mu_{k,n}$.  Each type has a physical deadline
$T_n$ and a cost rate $c_n$: a waiting job that exceeds its deadline
incurs cost.  The objective is the long-run average cost.

\paragraph{Uniformised MDP.}
The continuous-time dynamics are embedded into a discrete-time chain by
uniformisation.  At each \emph{period} exactly one event occurs:
an \textbf{arrival} of type $n$ (probability $\propto \lambda_n$),
a \textbf{service completion} (probability $\propto$ busy rates),
a \textbf{tick} that advances all waiting-time clocks by one unit
(probability $\propto \nu$, the tick rate\footnote{The RVI note denotes
the tick rate by $\gamma$; here we reserve $\gamma$ for the discount
factor and write $\nu$ for the tick rate.}), or
\textbf{nothing} (self-loop).  The uniformisation rate is
$\Lambda = \sum_n \lambda_n + \nu + \sum \text{(busy service rates)}$.
The state is $s = (f_0,\dots,f_{N-1},\, B,\, \alpha,\, \kappa)$: per-type
first-in-line (FIL) waiting times $f_n$ (in ticks; $-1$ if empty), the
busy configuration $B$, the action counter $\alpha$, and the category
$\kappa \in \{\textsc{AwaitEvent}, \textsc{AwaitAction}\}$.  Deadlines in
tick units are $d_n = T_n \nu$.

\paragraph{Action queue.}
Whenever idle capacity and waiting jobs coexist (after each tick), the
feasible (server, job-type) assignment candidates are collected into an
ordered \emph{action queue} (descending FIL, i.e.\ FIFO order) and
presented \emph{sequentially}.  At each position the agent chooses
action~$1$ (\emph{assign} the presented candidate) or action~$0$
(\emph{skip}, advance to the next candidate).  Skipping the final
candidate means the system idles until the next event.  Several decisions
may therefore occur within a single period, and a skipped job is
re-presented at the next tick.

\paragraph{Cost.}
The ground-truth objective (\texttt{reward\_type}~0) is the binary
deadline cost, charged at \textsc{AwaitEvent} states:
\begin{equation}
  c_{\mathrm{bin}}(s) = \sum_n c_n \,\mathbf{1}\{f_n > d_n\}.
  \label{eq:bincost}
\end{equation}
Policies are evaluated by simulation (100 trajectories $\times$
$5\cdot10^5$ periods) and reported as cost per unit time, i.e.\ mean cost
per period times $\Lambda$; we write this as
$\mathrm{cost}\times\Lambda$.  RVI supplies the exact optimal benchmark
$g^*$ on both instances, and the greedy FIFO policy (always assign the
oldest feasible candidate) is the natural heuristic baseline.

% ============================================================
\section{PPO in general terms}
\label{sec:ppo}
% ============================================================

\subsection{From value iteration to policy search}

RVI computes the exact optimal policy by sweeping the entire state space.
This is infeasible beyond small instances, so we turn to methods that
(i)~replace enumeration by \emph{simulation}, and (ii)~replace tables by
\emph{function approximation}.  Policy-gradient methods do both directly
in policy space: fix a parametric family of \emph{stochastic} policies
$\pi_\theta(a \mid s)$ --- here a neural network mapping a feature vector
of $s$ to a probability distribution over actions --- and improve
$\theta$ by stochastic gradient ascent on the expected return.

We use the discounted objective as a proxy for the average-cost
objective:
\begin{equation}
  J(\theta) = \E_{\pi_\theta}\Bigl[\sum_{t\ge 0} \gamma^{t}\, r_t\Bigr],
\end{equation}
with $\gamma$ close to $1$.  The effective horizon is
$1/(1-\gamma)$ periods, i.e.\ $1/\bigl((1-\gamma)\Lambda\bigr)$ time
units; with our $\gamma = 0.997$ this is ${\sim}330$ periods, long enough
to span many deadline windows.  (Rewards are negative costs; maximizing
return minimizes cost.)

\subsection{The policy gradient and the advantage function}

The policy gradient theorem states
\begin{equation}
  \nabla_\theta J(\theta)
  = \E_{\pi_\theta}\bigl[\nabla_\theta \log \pi_\theta(a_t \mid s_t)\,
    A^{\pi_\theta}(s_t, a_t)\bigr],
  \qquad
  A^{\pi}(s,a) = Q^{\pi}(s,a) - V^{\pi}(s).
  \label{eq:pg}
\end{equation}
The interpretation is direct: increase the log-probability of actions
that performed \emph{better than the policy's average} in their state
(positive advantage), decrease it otherwise.  Two consequences of
\eqref{eq:pg} matter throughout this note:
\begin{itemize}
  \item The update is \emph{advantage-weighted}: in states where all
        actions perform nearly equally under the current policy
        ($A \approx 0$), the gradient is negligible --- \emph{such states
        are never explicitly trained}, however frequently they are
        visited.
  \item The expectation is under the current policy's own state
        distribution: the data the algorithm learns from is generated by
        the very policy being trained, so a degraded policy also degrades
        its own future training data.
\end{itemize}

The advantage is unknown and must be estimated.  A \emph{critic}
$V_\phi(s)$ (a second network output) is regressed onto observed
returns, and the advantage is estimated by Generalized Advantage
Estimation (GAE): with the temporal-difference error
$\delta_t = r_t + \gamma V_\phi(s_{t+1}) - V_\phi(s_t)$,
\begin{equation}
  \hat A_t \;=\; \sum_{l \ge 0} (\gamma \lambda)^{l}\, \delta_{t+l},
  \label{eq:gae}
\end{equation}
where $\lambda \in [0,1]$ trades off bias ($\lambda = 0$: pure one-step
TD, biased by critic error) against variance ($\lambda = 1$: pure Monte
Carlo).\footnote{Throughout, plain $\lambda$ is the GAE parameter;
arrival rates always carry a subscript ($\lambda_n$).}

\subsection{The PPO clipped update}

Vanilla policy gradient permits exactly one gradient step per batch of
simulated experience: after the step, the data is off-policy and
\eqref{eq:pg} no longer holds.  PPO makes the data reusable for several
epochs of minibatch updates by controlling how far the new policy may
move from the policy $\pi_{\theta_{\mathrm{old}}}$ that generated the
data.  With the likelihood ratio
$\rho_t(\theta) = \pi_\theta(a_t \mid s_t)\,/\,
\pi_{\theta_{\mathrm{old}}}(a_t \mid s_t)$,
PPO maximizes the \emph{clipped surrogate}
\begin{equation}
  L^{\mathrm{clip}}(\theta) =
  \E_t\Bigl[\min\bigl(\rho_t \hat A_t,\;
      \clip(\rho_t,\, 1-\epsilon,\, 1+\epsilon)\, \hat A_t\bigr)\Bigr],
  \label{eq:clip}
\end{equation}
with $\epsilon = 0.2$.  The clip removes the incentive to push
$\rho_t$ beyond $1 \pm \epsilon$: once an action's probability has moved
by that factor, further movement earns no additional surrogate reward.
This is a cheap, first-order stand-in for a trust region.

The full loss combines three terms:
\begin{equation}
  L(\theta,\phi) \;=\; -L^{\mathrm{clip}}(\theta)
  \;+\; c_v \,\E_t\bigl[(V_\phi(s_t) - \hat R_t)^2\bigr]
  \;-\; c_e \,\E_t\bigl[\mathcal H\bigl(\pi_\theta(\cdot \mid s_t)\bigr)\bigr],
  \label{eq:loss}
\end{equation}
where $\hat R_t = \hat A_t + V_\phi(s_t)$ is the return target,
$\mathcal H$ is the policy entropy (a bonus that sustains exploration by
penalizing premature determinism), and $c_v, c_e$ are weights.  Both
heads share one network trunk, so the two regression problems interact
--- a fact that becomes important in Section~\ref{sec:vnorm}.

One training \emph{update} is: simulate $E$ parallel environments for
$T$ decisions each under the current policy; compute
\eqref{eq:gae} on the batch of $E \cdot T$ transitions; then take
several epochs of minibatch gradient steps on \eqref{eq:loss}.

\subsection{From stochastic policy to deterministic decision rule}
\label{sec:extraction}

PPO trains a stochastic policy, but what we ultimately want is a
deterministic scheduling rule.  The standard extraction is the
\emph{argmax policy}: in each state, take the action with the highest
probability.  It is worth flagging already here that this step is
\emph{not} innocuous: training only guarantees that the stochastic
policy performs well, and the argmax can differ from the stochastic
policy precisely in the states where the learned probabilities are close
to $1/2$.  Section~\ref{sec:forgiveness} shows that in our MDP this gap
is systematic rather than incidental, and
Section~\ref{sec:ablation} exhibits it empirically.

% ============================================================
\section{Our PPO formulation and the stabilization recipe}
\label{sec:ours}
% ============================================================

Applied naively, PPO fails on our MDP: on Experiment~2 it collapsed to
the ``never serve'' policy on 8 of 8 seeds ($9.9\times$ the optimal
cost).  Every ingredient below is the fix for a specific, diagnosed
failure mechanism; the ablation study in Section~\ref{sec:ablation}
substantiates each one in a controlled way.  We describe the
formulation mathematically; the concrete hyperparameters are collected
in Table~\ref{tab:hyper}.

\subsection{Semi-MDP timing: discounting by elapsed periods}
\label{sec:dper}

Decisions occur on the action queue, so consecutive decisions may be
separated by \emph{zero} periods (two candidates presented within the
same tick) or by \emph{many} periods (idle stretches between ticks).
The MDP between decisions is a semi-Markov process, and the discount
between decision $t$ and decision $t{+}1$ must be
\begin{equation}
  \gamma^{\Delta_t}, \qquad
  \Delta_t = \#\{\text{periods elapsed between the two decisions}\}
  \in \{0, 1, 2, \dots\},
\end{equation}
used both in the TD errors
$\delta_t = r_t + \gamma^{\Delta_t} V_\phi(s_{t+1}) - V_\phi(s_t)$
and in the GAE recursion (decay $\gamma^{\Delta_t}\lambda$).

\begin{mechanism}[Idleness subsidy through mis-discounting]
\label{mech:dper}
An earlier implementation discounted \emph{every} decision by one period
($\gamma^{\max(1,\Delta_t)}$, effectively).  An intra-tick decision
chain then discounts the future as if time passed at every skip, so a
policy that skips through the whole queue reaches the next (costly)
event with extra discount applied --- idleness is \emph{subsidized}, and
most strongly in busy states where the action queue is long.  This
single error was sufficient to push training into the never-serve
attractor.  The fix is the correct semi-MDP exponent with
$\gamma^0 = 1$.
\end{mechanism}

\subsection{Potential-based reward shaping (the urgency ramp)}
\label{sec:shaping}

The binary deadline cost \eqref{eq:bincost} is nearly flat: a job
incurs no cost until its deadline passes.  During early training almost
all sampled decisions therefore have $A \approx 0$, while the rare
deadline penalty arrives many decisions after the skip that caused it.
Against this sparse signal stands a strong attractor: never serving
makes the immediate future of every individual decision look identical,
and by Mechanism~\ref{mech:dper}-type credit-assignment noise the policy
drifts there easily.

We therefore shape the reward with a potential function.  Let
$\Phi(s)$ be the total \emph{accumulated urgency} of the jobs currently
waiting: each tick, the FIL job of type $n$ adds
\begin{equation}
  c_n \,\frac{\min(f_n, d_n)}{d_n}
  \label{eq:ramp}
\end{equation}
to its accumulated urgency (a ramp that rises linearly toward the
deadline, then saturates), and a job's accumulated urgency is
\emph{refunded} in full at the moment it enters service.  The shaped
per-transition cost is the base cost plus the accrual minus the refund,
i.e.\ exactly the potential-based form
$r' = r + \Phi(s) - \Phi(s')$ (in cost convention).  Potential-based
shaping leaves the optimal policy unchanged and, for stationary policies
under the average-cost criterion, leaves the long-run average unchanged
(the potential term telescopes).  We verified this empirically: FIFO and
RVI evaluate identically to four decimal places with and without
shaping, so RVI remains a valid benchmark for shaped training runs.

The effect on learning is dense, immediate credit: every tick spent not
serving an urgent job is charged \emph{now}, and serving refunds
\emph{now}.  Shaping is load-bearing: without it the success rate drops
from $6$--$7/8$ seeds to $2/8$ (Section~\ref{sec:ablation}).

\subsection{Value-target normalization}
\label{sec:vnorm}

Returns in this MDP are large (orders of $10^2$--$10^3$ in cost units),
while the policy-gradient part of \eqref{eq:loss} operates on
normalized advantages of order~$1$.

\begin{mechanism}[Value-scale domination]
With raw return targets, the value MSE term in \eqref{eq:loss} is orders
of magnitude larger than the policy term, and its gradients monopolize
the shared trunk: the representation is shaped for value regression
while the policy head starves.
\end{mechanism}

The fix: maintain a running estimate $\sigma_R$ of the return standard
deviation, train the value head on $\hat R_t / \sigma_R$, and rescale
its output by $\sigma_R$ wherever a raw value is needed (bootstrap
targets, GAE).  Advantages are additionally standardized per batch.
Without this ingredient, \emph{zero} of eight seeds succeed: every run
converges to one of the two attractor policies (never-serve or a FIFO
clone; Section~\ref{sec:ablation}).

\subsection{Persistent environments and staggered resets}
\label{sec:resets}

The $E = 16$ parallel simulation environments are \emph{persistent}:
the infinite-horizon chain continues across updates, which matches the
average-cost objective.  It also creates a trap.

\begin{mechanism}[Data poisoning by deep-queue excursions]
\label{mech:poison}
One bad excursion (e.g.\ a temporarily idle-happy policy) drives an
environment into a deep-backlog region where the shaping ramps
\eqref{eq:ramp} are saturated and no dense gradient exists.  If all
environments drift there, the training data no longer contains the
healthy short-queue states at all; the policy's behavior on those
states then degrades unmonitored (weights move on shared features), and
evaluation --- which starts from the empty state --- collapses.  This is
self-reinforcing through the on-policy data distribution and is not
recoverable by any gradient-side guard: only fresh-start data helps.
\end{mechanism}

The fix is \emph{staggered resets}: each update, one environment (round
robin, period $16$ updates) is re-initialized to the empty state.  The
data distribution stays anchored to the region the deployed policy
starts in, at a negligible cost in on-policyness.  Ablation: with no
resets, $1/8$ seeds succeed and five of eight converge to the exact
never-serve policy.

\subsection{Exploration and extraction: entropy and guarded
temperature annealing}
\label{sec:anneal}

\subsubsection{The re-presentation forgiveness mechanism}
\label{sec:forgiveness}

Recall from Section~\ref{sec:extraction} that the deployed policy is the
argmax of the trained stochastic policy.  In this MDP the gap between
the two is systematic.  A job skipped now is re-presented at every
subsequent tick while it waits.  Suppose the policy assigns probability
$p$ to serving a given job, and the job will be re-presented $m$ more
times within its deadline slack.  The probability it is served in time
is
\begin{equation}
  1 - (1-p)^{m+1},
\end{equation}
which for moderate $p$ and a few ticks of slack is already close to
$1$: with $p = 0.4$ and $m = 8$, the job is served in time with
probability ${\approx}0.99$.  \emph{The stochastic policy is
behaviorally near-optimal in such states even when $p$ is on the wrong
side of $\tfrac12$} --- so the advantage signal vanishes
(cf.\ the advantage-weighting in \eqref{eq:pg}) and training never
corrects the sign.  The argmax, however, turns $p = 0.4$ into
\emph{never serve this job here}, deterministically, at every
re-presentation.  Wrong-sign near-ties are invisible to training and
fatal to extraction.

\subsubsection{Guarded behavior-temperature annealing}

The mitigation is to make the \emph{behavior} policy converge toward the
argmax policy during training, so that states where the argmax action is
wrong actually hurt the observed return and receive gradient.  We sample
rollout actions from a tempered policy
\begin{equation}
  \pi^{(T)}_\theta(a \mid s) \;\propto\;
  \exp\bigl(z_\theta(s,a)/T\bigr),
\end{equation}
where $z_\theta$ are the policy logits and the temperature $T$ anneals
from $1$ toward $T_{\min}$ during the second half of training.
Annealing \emph{blindly} on a clock is dangerous --- it amplifies
whichever mode the policy is currently in, sharpening bad seeds into
collapse --- so the schedule is \emph{guarded}: $T$ steps down (factor
$0.9$) only while an exponential moving average of the rollout reward
rate is not degrading relative to its best level since annealing began,
and steps back up when it is.  The health signal is measured
\emph{per period}, never per decision: an idling policy makes many
zero-reward intra-tick decisions, so a per-decision mean dilutes
precisely when the policy degrades toward idleness.

Two auxiliary devices complete the scheme.  A \emph{best-sharp snapshot}
stores the network parameters whenever behavior is both sharp
($T \le 2T_{\min}$) and healthy, and is restored at the end of training
unless the final parameters match it --- the guard's late retreats
otherwise discard the best policy of the run.  At evaluation time we
extract a small \emph{readout battery} (argmax; tempered stochastic;
argmax with a bias added to the serve logit) and select by evaluation;
on hard instances this converts a per-seed success rate into a
per-problem one.  Separately, the entropy coefficient $c_e$ in
\eqref{eq:loss} is annealed linearly to zero over the second half of
training, so the policy is allowed to become deterministic exactly when
the temperature schedule wants it to.

\subsection{What remains: advantage-weighted blindness}
\label{sec:blindness}

The mechanisms above remove the systematic failures; what remains on the
hardest instance is a residual \emph{seed lottery}
(Section~\ref{sec:exp2}).  Its root cause is the advantage-weighting in
\eqref{eq:pg} combined with Section~\ref{sec:forgiveness}: some
decision boundaries are behaviorally flat under the stochastic policy
and never receive gradient, and whether the argmax lands on the right
side of them is decided by initialization noise.  More compute helps
some seeds (doubling to $30{,}000$ updates rescues several
$6{,}000$-update failures); an auxiliary advantage head (a third network
output regressing the observed GAE advantages, giving an alternative
readout) rescued individual seeds but showed no systematic improvement
in an 8-seed A/B test ($5/8$ vs $5/8$).  The principled fix --- an
action-space reformulation in which each freed server triggers exactly
one decision (``which class to serve, or idle'') so that skipping is a
real commitment rather than a re-presentable preference --- is known but
deliberately deferred; a first probe in that direction is reported in
Section~\ref{sec:skipall}.

\subsection{Hyperparameters}

\begin{table}[h]
\centering
\caption{The tuned recipe (all experiments in
Section~\ref{sec:experiments} unless stated otherwise).}
\label{tab:hyper}
\begin{tabular}{llll}
\toprule
discount $\gamma$ & $0.997$ &
environments $E$ & $16$ (persistent) \\
GAE $\lambda$ & $0.95$ &
rollout $T$ & $512$ decisions \\
clip $\epsilon$ & $0.2$ &
updates & $6000$ \\
entropy coef.\ $c_e$ & $0.03$, annealed $\to 0$ &
epochs $\times$ minibatch & $4 \times 256$ \\
value coef.\ $c_v$ & $0.5$ &
learning rate & $3\cdot10^{-4}$ (Adam) \\
temperature floor $T_{\min}$ & $0.1$, guarded &
reset period & $16$ updates \\
network & \multicolumn{3}{l}{MLP $64 \times 32$, shared trunk; policy,
value (and aux-advantage) heads} \\
\bottomrule
\end{tabular}
\end{table}

Two calibration notes.  First, $\gamma = 0.997$ is load-bearing:
$0.99$ plateaus at a FIFO clone (effective horizon too short to see the
benefit of strategic idling), $0.999$ destabilizes credit assignment;
the sensitivity sweep (three seeds per knob) puts both alternatives at
$3$--$5\times$ higher mean cost ratios.  Second, the variance knobs
matter: raising the entropy coefficient from $0.01$ to $0.03$ and
lowering $T_{\min}$ from $0.25$ to $0.1$ each reduced seed-to-seed
standard deviation by a factor ${\sim}4$ at equal mean.

% ============================================================
\section{Experiments and results}
\label{sec:experiments}
% ============================================================

\subsection{The two benchmark instances}

Both instances use tick rate $\nu = 3$ and the binary deadline cost
\eqref{eq:bincost} as ground truth (training uses its shaped version,
Section~\ref{sec:shaping}).  Reported numbers are cost per unit time
($\mathrm{cost}\times\Lambda$); lower is better.

\paragraph{Experiment 3 (routing: specialist + generalist).}
$N = 2$ types, $\lambda = (0.8, 0.2)$, costs $c = (100, 300)$, deadline
$T_n = 1$; server pool $0$ is a \emph{specialist} (type $0$ only,
$\mu = 1$), pool $1$ a \emph{generalist} (both types, $\mu = 1$).  The
optimal policy must protect the generalist's capacity for type-$1$ jobs.
Benchmarks: $\mathrm{RVI}\times\Lambda = 18.96$,
$\mathrm{FIFO}\times\Lambda = 24.15$ (ratio $1.27$).

\paragraph{Experiment 2 (strategic idling: fully flexible, asymmetric
cost).}
$N = 2$ types, $\lambda = (0.25, 0.25)$, costs $c = (100, 300)$,
deadlines $T_n = 3$; two fully flexible single-server pools with
$\mu = 0.35$ ($\rho \approx 0.71$).  Since every server can serve every
job, the \emph{only} way to beat FIFO is strategic idling: refusing a
cheap type-$0$ job now to keep capacity free for a possible expensive
type-$1$ arrival.  Benchmarks: $\mathrm{RVI}\times\Lambda = 40.50$,
$\mathrm{FIFO}\times\Lambda = 70.46$ (ratio $1.74$).  This is the harder
instance: the payoff of idling is realized only through rare future
events, exactly the signal that Sections~\ref{sec:forgiveness}
and~\ref{sec:blindness} identify as hardest to learn.

``Within 5\%'' below means argmax cost ratio to RVI at most $1.05$.

\subsection{Experiment 3: solved}
\label{sec:exp3}

With the recipe of Section~\ref{sec:ours} (at the pre-tuning variance
settings $c_e = 0.01$, $T_{\min} = 0.25$, rollout $256$, plus the
auxiliary head), the eight seeds give
\[
  \mathrm{argmax}\times\Lambda =
  19.12,\; 19.88,\; 19.16,\; 18.97,\; 19.37,\; 18.95,\; 19.09,\; 21.55
\]
against $\mathrm{RVI}\times\Lambda = 18.96$: \textbf{7/8 within 5\%}
(the eighth at $13.7\%$), all far below FIFO's $24.15$.  Training is
${\sim}10$ minutes per seed on a shared 16-core node slice.  We consider
this instance solved and use it below mainly as the easy control.

\subsection{Experiment 2: solved per problem, seed lottery per run}
\label{sec:exp2}

Table~\ref{tab:exp2} shows the argmax results across eight seeds for
three budgets of the original recipe, and for the variance-tuned recipe
(Table~\ref{tab:hyper}).

\begin{table}[h]
\centering
\caption{Experiment 2, argmax $\mathrm{cost}\times\Lambda$ vs
$\mathrm{RVI}\times\Lambda = 40.50$ (FIFO $= 70.46$).  $\times$ denotes
collapse ($> 2\times$ RVI).  ``+aux'' adds the auxiliary advantage head;
``tuned'' is Table~\ref{tab:hyper}.}
\label{tab:exp2}
\begin{tabular}{lcccc}
\toprule
Seed & $6000$u & $6000$u+aux & $30\,000$u & tuned, $6000$u \\
\midrule
1 & 44.78 & 40.80 & 40.85 & 42.84 \\
2 & $\times$ & 41.99 & 41.73 & 42.20 \\
3 & 44.81 & 53.27 & 60.84 & 40.78 \\
4 & $\times$ & 42.00 & 41.58 & 42.51 \\
5 & 40.81 & $\times$ & 40.78 & 41.02 \\
6 & 40.92 & 41.59 & 42.37 & 40.78 \\
7 & 40.90 & 40.89 & 41.89 & 40.93 \\
8 & $\times$ & 57.48 & $\times$ & $\times$ \\
\midrule
within 5\% & 4/8 & 5/8 & 6/8 & 7/8\footnotemark \\
\bottomrule
\end{tabular}
\end{table}
\footnotetext{Seeds 1 and 4 sit at $5.8\%$ and $5.0\%$; counting
strictly, 5/8 within 5\% and 7/8 within 6\%.}

Three observations.  First, the tuned recipe is both better and much
tighter: seven of eight seeds land in $[40.78,\, 42.84]$.  Second,
failures do not attach to specific seeds: seed 3 fails with the
auxiliary head and succeeds without; seed 5 the reverse; only seed 8
fails in most (not all) configurations.  This is a lottery over
initializations, not a set of hard seeds --- consistent with
Section~\ref{sec:blindness}.  Third, and practically decisive: training
2--3 seeds and selecting by evaluation yields a within-5\% policy
\emph{every} time; per problem instance, Experiment~2 is solved at a
cost of ${\sim}30$ minutes of compute.  The runs are exactly
reproducible (same seed, same result to four decimals), so the tables
above are deterministic artifacts, not noisy samples.

\subsection{Controlled ablations: substantiating each ingredient}
\label{sec:ablation}

The mechanisms of Section~\ref{sec:ours} were diagnosed during
development, i.e.\ on a moving codebase.  To substantiate them, we
re-ran everything as controlled single-ingredient ablations from the
frozen tuned recipe: eight cells $\times$ eight seeds, identical in
every respect except the named ingredient
(Table~\ref{tab:ablation}).

\begin{table}[h]
\centering
\caption{Ablation matrix on Experiment 2: argmax cost ratio to RVI,
8 seeds per cell.  Every cell is the tuned recipe with exactly one
change.  For reference: never-serve evaluates at ratio $9.876$, FIFO at
$1.738$.}
\label{tab:ablation}
\begin{tabular}{llccc}
\toprule
Cell & Removes / restores & within 5\% & median & worst \\
\midrule
\texttt{base} & (control) & 6/8 & 1.03 & 5.41 \\
\texttt{dper} & restores mis-discounting (M.\ref{mech:dper}) & 0/8 & 7.26 & 9.85 \\
\texttt{vnorm} & value normalization off (\S\ref{sec:vnorm}) & 0/8 & 9.88 & 9.88 \\
\texttt{resets0} & no environment resets (\S\ref{sec:resets}) & 1/8 & 9.88 & 9.88 \\
\texttt{shape0} & no reward shaping (\S\ref{sec:shaping}) & 2/8 & 1.87 & 9.69 \\
\texttt{anneal0} & no temperature annealing (\S\ref{sec:anneal}) & 3/8 & 1.29 & 7.75 \\
\bottomrule
\end{tabular}
\end{table}

The matrix confirms the recipe ingredient by ingredient, and the
\emph{failure signatures} carry as much information as the counts:

\begin{itemize}
  \item \textbf{Mis-discounting} (\texttt{dper}): total failure, with
  most seeds at ratios $8.9$--$9.85$, i.e.\ at or near the never-serve
  attractor --- the direction Mechanism~\ref{mech:dper} predicts
  (idleness subsidy), not merely ``worse''.
  \item \textbf{Value scale} (\texttt{vnorm}): perfectly bimodal.  Every
  seed converges to \emph{exactly} ratio $9.876$ (never-serve; the value
  matches an independently simulated always-skip policy to four
  decimals) or \emph{exactly} $1.738$ (a FIFO clone, matching simulated
  FIFO to four decimals).  With a crushed trunk the policy can only fall
  into one of the two attractors.
  \item \textbf{Resets} (\texttt{resets0}): seven of eight seeds
  collapse (five to exact never-serve); one lucky seed solves the
  instance --- Mechanism~\ref{mech:poison} is about an excursion
  \emph{trap}, and a seed that never has the excursion never falls in.
  \item \textbf{Shaping} (\texttt{shape0}): success drops to 2/8 with
  three full never-serve collapses.  The honest nuance: shaping is not
  binary --- occasional seeds succeed on the sparse cost --- but it
  roughly triples the success probability.
  \item \textbf{Temperature annealing} (\texttt{anneal0}): the most
  instructive cell.  Failures are mostly graceful (ratios
  $1.28$--$2.46$), and on four of the five failed seeds the
  \emph{stochastic} readout is nearly optimal while the argmax is far
  off, e.g.\ (argmax vs stochastic, $\times\Lambda$):
  $94.2$ vs $42.0$;\; $52.5$ vs $41.7$;\; $314.0$ vs $43.1$;\;
  $51.7$ vs $41.3$.
  Without annealing, PPO \emph{learns} the right behavior but the
  argmax cannot \emph{extract} it --- the re-presentation forgiveness
  gap of Section~\ref{sec:forgiveness}, measured directly.
\end{itemize}

\subsection{A first action-space probe: the skip-all action}
\label{sec:skipall}

As a low-cost probe toward the reformulation of
Section~\ref{sec:blindness}, we added an optional third action:
\emph{skip-all} declines every remaining candidate in the current
action queue in one decision, making whole-tick idling a single,
readable choice instead of a chain of skips.  By construction skip-all
is value-degenerate with the skip chain (verified: a hand-coded
always-skip-all policy and an always-skip policy evaluate identically to
four decimals), so RVI benchmarks are unaffected; the policy head simply
gains a third logit, initialized pessimistically (bias $-3$,
${\approx}2\%$ initial probability) so that uniform initialization does
not grant whole-tick idling a $1/3$ prior.

The 8-seed result is negative but mechanistically clean: 3/8 seeds
within 5\% (vs 6/8 for the control), with four seeds collapsing --- and
on the collapsed seeds the stochastic readout remains moderate
($\times\Lambda$ between $51$ and $92$) while the argmax is
catastrophic ($352$--$372$).  Skip-all does not break \emph{learning};
it breaks \emph{extraction}: a wrong-sign near-tie on the skip-all
logit, harmless under the stochastic policy, is repeated
deterministically at every tick under the argmax and idles the system.
This sharpens the design requirement for the full reformulation: it is
not enough to make idling expressible in one action --- skipping must be
a \emph{commitment} (no re-presentation within the tick) so that
training itself feels the cost of wrong idling decisions, rather than
deferring the problem to extraction.

\subsection{Summary}

\begin{itemize}
  \item PPO, correctly adapted to the semi-MDP structure and stabilized
  by four ingredients (correct elapsed-period discounting,
  potential-based urgency shaping, value-target normalization, staggered
  resets) plus guarded temperature annealing for extraction, learns
  RVI-optimal policies on both benchmark instances.
  \item Experiment 3 (routing) is solved: 7--8/8 seeds within 5\% of
  RVI.  Experiment 2 (strategic idling) is solved per problem instance:
  7/8 seeds within 6\% under the tuned recipe, and 2--3 seeds with
  evaluation-based selection always suffice.
  \item Each ingredient is substantiated by a controlled 8-seed
  ablation; removing any one collapses success to 0--3/8, with failures
  landing on two exactly reproducible attractors (never-serve, ratio
  $9.876$; FIFO clone, ratio $1.738$).
  \item The residual seed lottery on Experiment 2 traces to
  advantage-weighted blindness at behaviorally flat decision boundaries,
  aggravated by re-presentation forgiveness at extraction time.  The
  skip-all probe shows that naive action-space patches inherit the
  extraction problem; the per-event reformulation, which removes
  re-presentation by construction, is the principled next step.
\end{itemize}

\end{document}
